\section*{第 \ref{ch:g10:biodiversity-scales} 章——各尺度上的生物多样性}

\begin{solution}{exo:g10:biodiversity-scales:1}
\omterm{def:g10:biodiversity-scales:biodiversity}{生态系统}（池塘、山毛榉林）；\omterm{def:g10:biodiversity-scales:species}{物种}（池塘中的青蛙、
蝾螈和芦苇）；\omterm{def:g10:universal-dna:gene}{基因}（人类
种群中 ABO \omterm{def:g10:universal-dna:gene}{基因}的\omterm{def:g10:universal-dna:gene}{等位基因}）。
\end{solution}

\begin{solution}{exo:g10:biodiversity-scales:2}
相似生物体组成的群体，它们能够相互交配，并产生可育的
后代。狮与虎产生的是不育的杂交后代，因此它们是
两个\omterm{def:g10:biodiversity-scales:species}{物种}。
\end{solution}

\begin{solution}{exo:g10:biodiversity-scales:3}
$\num{1000000}/6400 \approx 156$：约多 150 倍。
\end{solution}

\begin{solution}{exo:g10:biodiversity-scales:4}
种群中某\omterm{def:g10:universal-dna:gene}{基因}全部拷贝中属于给定
\omterm{def:g10:universal-dna:gene}{等位基因}的比例。B $= 1 - 0.66 - 0.27 = 0.07$。
\end{solution}

\begin{solution}{exo:g10:biodiversity-scales:5}
共五次；其中最大的是二叠纪末\omterm{prop:g10:biodiversity-scales:changes}{大灭绝}，发生在
2.52 亿年前，使约 95\% 的海洋\omterm{def:g10:biodiversity-scales:species}{物种}灭绝。
\end{solution}

\begin{solution}{exo:g10:biodiversity-scales:6}
第二个：\omterm{def:g10:biodiversity-scales:species}{物种}丰富度相同，但相对丰度分布均匀得多。
仅凭丰富度无法区分二者，相对丰度才能区分。
\end{solution}

\begin{solution}{exo:g10:biodiversity-scales:7}
两块样地相同，只是同学的取样量多了十倍。\omterm{def:g10:biodiversity-scales:species}{物种}
数随取样量增加而增多，直到曲线趋于平缓，因此取样量不同时
所得计数不可比较——该生记录的 8 种是取样不足造成的，
而非样地更为贫瘠。
\end{solution}

\begin{solution}{exo:g10:biodiversity-scales:8}
少数幸存个体只携带了该\omterm{def:g10:biodiversity-scales:species}{物种}\omterm{def:g10:universal-dna:gene}{等位基因}的一部分；
它们的后代无论数量多少，都只拥有这一缩减后的组合。\omterm{def:g10:biodiversity-scales:biodiversity}{遗传
多样性}在瓶颈效应中丧失，不会因数量增加而恢复：
该\omterm{def:g10:biodiversity-scales:species}{物种}面对新疾病或气候变化时别无储备，
猎豹即是一例。
\end{solution}

\begin{solution}{exo:g10:biodiversity-scales:9}
灭绝前约 560 个科，灭绝后约 300 个，
科的损失比例为 $260/560 \approx 45\%$。只有一个科的全部\omterm{def:g10:biodiversity-scales:species}{物种}都灭绝时，该科才会消失，因此
许多科仅凭一两个\omterm{def:g10:biodiversity-scales:species}{物种}即得以存留，而其余
\omterm{def:g10:biodiversity-scales:species}{物种}已经灭绝；\omterm{def:g10:biodiversity-scales:species}{物种}损失的比例（约 95\%）远大于
科损失的比例。
\end{solution}

\begin{solution}{exo:g10:biodiversity-scales:10}
每个种群都携带全部三种\omterm{def:g10:universal-dna:gene}{等位基因}，其成员
之间只是携带的\omterm{def:g10:universal-dna:gene}{等位基因}组合不同；种群之间的
差异仅仅是频率的变化。因此，该\omterm{def:g10:biodiversity-scales:species}{物种}的大部分\omterm{def:g10:biodiversity-scales:biodiversity}{遗传
多样性}存在于各个种群内部：两个相邻个体
之间的差异位点数，往往与两个相距甚远的人一样多。
\end{solution}

\begin{solution}{exo:g10:biodiversity-scales:11}
\omterm{def:g10:biodiversity-scales:biodiversity}{生态系统}：池塘本身及其群落。\omterm{def:g10:biodiversity-scales:species}{物种}：仅局部生活在那里的
蝾螈或蜻蜓。\omterm{def:g10:universal-dna:gene}{基因}：池塘中青蛙种群所携带的
\omterm{def:g10:universal-dna:gene}{等位基因}，可能与
其他池塘中的不同。
\end{solution}

\begin{solution}{exo:g10:biodiversity-scales:12}
若\omterm{def:g10:biodiversity-scales:species}{物种}的存续期约为 $\num{2e6}$ 年，则 $\num{2e6}$ 个
\omterm{def:g10:biodiversity-scales:species}{物种}平均每年约灭绝一个。若速度快上一千倍，则每年约 1000
个\omterm{def:g10:biodiversity-scales:species}{物种}，即每世纪 $\num{100000}$ 个，约占已知\omterm{def:g10:biodiversity-scales:species}{物种}的 5\%
在一百年内。
\end{solution}

\begin{solution}{exo:g10:biodiversity-scales:13}
这次灭绝清空了恐龙曾占据的多种生活方式——大型
植食动物、大型捕食者、飞行动物和游泳动物。
幸存的哺乳动物谱系摆脱了这些竞争者，并凭借
自身的\omterm{def:g10:biodiversity-scales:biodiversity}{遗传多样性}，适应了这些空缺的生态位，
分化出新的\omterm{def:g10:biodiversity-scales:species}{物种}：这次灭绝开辟了机会，
新的多样性填补了这些空白。
\end{solution}

\begin{solution}{exo:g10:biodiversity-scales:14}
栽培小麦已丢失了其祖先的大部分\omterm{def:g10:universal-dna:gene}{等位基因}。野生
近缘种仍保留着这些\omterm{def:g10:universal-dna:gene}{等位基因}：抗某种病害、
抗干旱、耐盐的\omterm{def:g10:universal-dna:gene}{等位基因}，是当前品种所没有的。将其导入
小麦即可恢复这一选择；一旦某种野生近缘种灭绝，其\omterm{def:g10:universal-dna:gene}{等位基因}
便永久丧失。
\end{solution}

\begin{solution}{exo:g10:biodiversity-scales:15}
灭绝本是自然现象，但自然速率约为每年一种；当前
的速率快了数百至一千倍，与\omterm{prop:g10:biodiversity-scales:changes}{大灭绝}相当，
而从\omterm{prop:g10:biodiversity-scales:changes}{大灭绝}中恢复所需的时间长达数百万年，
超过了我们这一\omterm{def:g10:biodiversity-scales:species}{物种}存在的时间。与此同时，正在消失的\omterm{def:g10:biodiversity-scales:biodiversity}{生态系统}
为我们提供食物、传粉、土壤、水净化以及作物的
遗传资源。这一先例表明，损失是真实的，而
在人类可感知的时间尺度上，恢复是无法实现的。
\end{solution}

\begin{solution}{pb:g10:biodiversity-scales:1}
\textbf{1.} \omterm{def:g10:biodiversity-scales:species}{物种}数分别为 31 和 9：比值 $31/9 \approx 3.4$。

\textbf{2.} 农场 A 中没有\omterm{def:g10:biodiversity-scales:species}{物种}的占比超过 22\%，而农场
B 上有一个\omterm{def:g10:biodiversity-scales:species}{物种}占 84\%。均匀度说明个体在所发现的各\omterm{def:g10:biodiversity-scales:species}{物种}之间如何分配；
一个丰富但不均匀的群落由单一\omterm{def:g10:biodiversity-scales:species}{物种}占主导，
其余\omterm{def:g10:biodiversity-scales:species}{物种}都很稀少。

\textbf{3.} 增量逐渐减小（7、4、2）：曲线趋于平缓，
样本已接近完整；再多取几个样方，最多
只能增加一两个\omterm{def:g10:biodiversity-scales:species}{物种}。

\textbf{4.} 农场 B 的曲线在第 15 个样方时已达到 9，之后再未增加：
完整的取样曲线会趋于平台，
农场 B 已达到这一平台，农场 A 则接近达到。

\textbf{5.} 许多植物只有开花时才能看见或辨认，
昆虫和鸟类也随季节变化；换一个月份取样，
数目的变化可能与树篱无关。

\textbf{6.} 昆虫 $46/14 \approx 3.3$；鸟类 $17/5 = 3.4$。
三个类群的比值均约为 3，结论一致。

\textbf{7.} 树篱提供昆虫取食和繁殖所依赖的植物（叶、花、枯木、
遮蔽处）；植物\omterm{def:g10:biodiversity-scales:species}{物种}减少，意味着昆虫\omterm{def:g10:biodiversity-scales:species}{物种}也减少；
昆虫以及树篱上的浆果和巢址为鸟类提供食物和栖身之所，
因此鸟类多样性也随之下降。

\textbf{8.} 不能：产量只反映作物本身。树篱可能为田地遮挡风力、
保持土壤，并为传粉者和作物害虫的天敌提供栖身之所，
其作用可能要在歉收年份或数十年间才显现，
而非体现在某一季的产量数字中。

\textbf{9.} 在树篱中筑巢的\omterm{def:g10:biodiversity-scales:species}{物种}，如黑鸫或篱雀
（或树篱蜗牛、山楂盾蝽）：它们需要灌木用于筑巢、
取食或栖身，而农场 B 已不再提供这些。

\textbf{10.} \omterm{def:g10:biodiversity-scales:biodiversity}{生态系统}层面：树篱本身作为一种栖息地被移除，
\omterm{def:g10:biodiversity-scales:species}{物种}和\omterm{def:g10:universal-dna:gene}{基因}随之消失；
其他层次上的损失都是后果。

\textbf{11.} 农场 A：y 的拷贝数为 $2 \times 100 + 200 = 400$，b 的拷贝数
为 $200 + 2 \times 100 = 400$，共 800 个，y 的频率为 0.5。农场 B：y
$8 + 32 = 40$，b 为 $32 + 128 = 160$，共 200 个，y 的频率为 0.2。

\textbf{12.} 农场 A：两种\omterm{def:g10:universal-dna:gene}{等位基因}频率相等，50\% 的蜗牛
同时携带两者；农场 B：32\% 的蜗牛同时携带两者，且 b 占优势。农场 A 的
多样性更高。

\textbf{13.} 在裸露的河岸上，黄色蜗牛更容易被捕食；
一代又一代，传递下去的 y 拷贝数越来越少，
y 的频率由约 0.5 降至 0.2。

\textbf{14.} 优势会发生逆转：黄色蜗牛在有遮荫的树篱中更易存活，
y 的频率将在世代更替中再次上升
——前提是该\omterm{def:g10:universal-dna:gene}{等位基因}仍然存在。

\textbf{15.} 农场 B 的情况：在只有几百只的种群中，一个糟糕的年份
或少数偶然死亡，就可能使较稀有的\omterm{def:g10:universal-dna:gene}{等位基因}的全部拷贝消失；
在数万只的种群中，该\omterm{def:g10:universal-dna:gene}{等位基因}由数千只蜗牛携带，
偶然性无法将其抹去。

\textbf{16.} 农场 A：$31 + 46 + 17 = 94$；农场 B：$9 + 14 + 5 = 28$；
比值 $94/28 \approx 3.4$。

\textbf{17.} 仅占土地 3\% 的部分，保存着
农场约三分之二\omterm{def:g10:biodiversity-scales:species}{物种}的栖息地。

\textbf{18.} 植物 $24/31 \approx 77\%$；昆虫 $35/46 \approx 76\%$；
鸟类 $12/17 \approx 71\%$。

\textbf{19.} \omterm{def:g10:biodiversity-scales:biodiversity}{遗传多样性}：\omterm{def:g10:biodiversity-scales:species}{物种}可从邻近种群中回归，
但本地种群中丢失的\omterm{def:g10:universal-dna:gene}{等位基因}（例如农场 B 的蜗牛所携带的）
只有在迁入者恰好携带它们时才会回来；
重新种植树篱并不能重新产生已丢失的\omterm{def:g10:universal-dna:gene}{等位基因}。

\textbf{20.} 树篱使农场各统计类群的\omterm{def:g10:biodiversity-scales:species}{物种}丰富度
约为原来的 3.4 倍；重新种植的树篱能恢复\omterm{def:g10:biodiversity-scales:biodiversity}{生态系统}，
并在一段时间后带回大多数\omterm{def:g10:biodiversity-scales:species}{物种}，
但无法独自恢复本地种群已丢失的\omterm{def:g10:universal-dna:gene}{等位基因}。
\end{solution}
